\noindent In many real-world scenarios, understanding whether an event will
occur is important. Nevertheless, understanding when it will occur
is even more informative. Consider, for example, a patient with a
chronic disease. Clinicians are often interested not only in the
effectiveness of a treatment, but also in how long the patient is
expected to survive under different therapeutic strategies. Similarly,
in the financial sector, institutions may aim to monitor the time until
a customer defaults on a loan, seeking to better manage risk exposure.
In this regard, companies are also frequently concerned with customer
churn, that is, the time until a client stops using a service or cancels
a subscription. Anticipating when a customer is likely to leave allows
organizations to implement targeted retention strategies and improve
long-term profitability. All of these examples are ubiquitous, and
survival analysis is particularly concerned with modeling the time
until an event of interest especialy in the presence of incomplete
observations.

\section{Characteristics of Survival Data}

According to \citeonline[p.~2]{moore2016applied}, survival data are defined by two
characteristics. The first one is that the response variable is a non-negative discrete
or continuous random variable, and represents the time from a well-defined origin to a
well-defined event. For example, the time elapsed between the starting point of a cancer
diagnosis and the occurrence of a specified event of interest (e.g., the pacient death).
The second on is censoring. Censoring occurs when the exact time of the event of interest
is not fully observed. When censoring does not occurs, it is said that the individual
failed.

Formally, let $T^*$ denote a random variable representing the time to failure and $U$ is
a random variable representing the time to a censoring event. Therefore, the observed
survival time is denoted as $T=\text{min}\left(T^*, U\right)$. Following the previous
definition,

$$
\delta = \begin{cases}
1 & \text{if } T^* < U \\
0 & \text{if } T^* \geq U
\end{cases}
$$

\noindent denotes a censoring indicator. That is,
$\delta$ is 1 if $T$ is a failured time and 0 if $T$ is a censored time.




% \citeonline[p.~2]{moore2016applied} defines survival data as two characteristics.
% The first one is that the response variable is a non-negative discrete or
% continuous random variable, and representes the time from a well-defined
% origin to a well-defined event (i.e., the survival time). The second one
% is censoring, is the most important. Censoring happens when the starting or
% ending events are not precisely observed. When the censoring does not happens,
% it means that the event was observed in the survival time and in this case is
% to be said a failure.



% In other
% words, it means that the main characteristic of survival data is
% that the response variable is a non-negative discrete or continuous
% random variable, and represents the time from a well-defined origin to a
% well-defined
